\scanpage{50}
\mspage{6}
% archive sheet number: 49
\uline{Remarque} En fait, on a un foncteur plus précis, vers une catégorie de graphes munis d'une structure où la partie $\rho$ de la pondération est \deleted{\unsure{complétée}} \inserted{\unsure{remplacée}} par une structure plus précise, à savoir \deleted{une \unsure{application}} \deleted{\unsure{une appli.}} famille
\[
\deleted{\unsure{A|\Delta}\,(\unsure{\widetilde R,\alpha,p:R\to\vec A}\ \underset{\vec A/\unsure{\sigma|\alpha}}{\Vert})}\quad (R_a)_{a\in A}
\]
% slanted
\margin{$\omega(a)\overset{\mathrm{def}}{=}\{\vec a\in\vec A\ /\ \vec a\text{ au dessus de }a\}$}
d'ensembles, et pour $\forall a\in A$, si $\omega(a)$ désigne l'ens des ``orientations'' de $a$ ($\unsure{\mathrm{card}\,\omega(a)|\mathrm{card}\,|\omega(a)|}\in\{1,2\}$), une application
\deleted{\unsure{un ensemble fini} $R$, et $p$ une application $R\longrightarrow A$ \unsure{composée}}
\diagram{
\begin{tikzcd}[ampersand replacement=\&]
\omega\longmapsto u_{\omega}:\quad\omega(a)\arrow[r] \& \mathrm{Circ}(R_{a})
\end{tikzcd}
}
ou des permutations circulaires sur $R_{a}$
\deleted{et \unsure{aussi} une \unsure{autre} application \uline{\unsure{injective}}}
\diagram{
\begin{tikzcd}[ampersand replacement=\&]
\deleted{\unsure{R/\alpha}}\arrow[r,hook,"\times"description]\&\deleted{\unsure{\vec A|\widetilde A}}
\end{tikzcd}
}
avec la condition que si $\omega,\omega'\in\omega(a)$, $\omega\ne\omega'$, alors $\deleted{u_{\omega}u_{\omega'}=\mathrm{id}_{R_{a}}}\,u_{\omega'}=u_{\omega}^{-1}$

On peut aussi considérer que la donnée supplémentaire revient à la donnée de
\diagram{
\begin{tikzcd}[ampersand replacement=\&]
\underset{\substack{\Vert\\a\mapsto\mathrm{Card}\,p^{-1}(a)}}{r:}\quad A\arrow[r]\&\mathbb N
\end{tikzcd}
}
(la ``pondération-arêtes''), et la donnée d'une famille de torseurs
\[
(R_{a})_{a\in A},\quad \deleted{\unsure{\text{avec}}}
\]
$R_a$ un torseur sous $\mathbb{Z}/r(a)\mathbb{Z}\wedge_{\{\pm1\}}^{\omega(a)}$. On associe à $(X,\widetilde{\Sigma},R)$ une telle structure supplémentaire
\[
R_{a}=R\cap\widetilde{\Sigma}_{a}\quad
(\widetilde{\Sigma}_{a}\ \text{la comp. connexe de }\widetilde{\Sigma}\
\deleted{\unsure{\text{correspondant}}}\ \text{qui correspond à } a\in A\simeq\pi_{0}(\widetilde{\Sigma})),
\]
\diagram{
\begin{tikzcd}[ampersand replacement=\&]
\deleted{\unsure{R=R(X_C)|R=R(\ast)}}\ ,\quad\deleted{R}\arrow[r,"\times"description]\arrow[d,equal]\&
\deleted{S}\arrow[d,equal]\\
\deleted{\unsure{\Gamma}}\arrow[r,"\times"description]\&\deleted{\pi_0(\widetilde{\Sigma})}
\end{tikzcd}
}
et $\omega\mapsto u_{\omega}$ défini de façon évidente.
\deleted{$p$ déduit de l'ordre circulaire sur les $R_a$}

\uline{NB} Même si $\unsure{r(\omega)|r(a)}=1$ ou $r(\omega)=2$, le fait de disposer
